\documentclass[UTF8, 12pt, fontset=fandol, oneside]{ctexart}
\usepackage{researchnotes}

\title{机器学习：提升方法}
\author{}
\date{}

\begin{document}

\maketitle

\section*{第 8 章\quad 提升方法}

提升（boosting）方法是一种常用的统计学习方法，应用广泛且有效。在分类问题中，它通过改变训练样本的权重，学习多个分类器，并将这些分类器进行线性组合，提高分类的性能。

本章首先介绍提升方法的思路和代表性的提升算法 AdaBoost；然后通过训练误差分析探讨 AdaBoost 为什么能够提高学习精度；并且从前向分步加法模型的角度解释 AdaBoost；最后叙述提升方法更具体的实例——提升树（boosting tree）。AdaBoost 算法是 1995 年由 Freund 和 Schapire 提出的，提升树是 2000 年由 Friedman 等人提出的。

\subsection*{8.1\quad 提升方法 AdaBoost 算法}

\subsection*{8.1.1\quad 提升方法的基本思路}

提升方法基于这样一种思想：对于一个复杂任务来说，将多个专家的判断进行适当的综合所得出的判断，要比其中任何一个专家单独的判断好。实际上，就是“三个臭皮匠顶个诸葛亮”的道理。

历史上，Kearns 和 Valiant 首先提出了“强可学习”（strongly learnable）和“弱可学习”（weakly learnable）的概念。指出：在概率近似正确（probably approximately correct，PAC）学习的框架中，一个概念（一个类），如果存在一个多项式的学习算法能够学习它，并且正确率很高，那么就称这个概念是强可学习的；一个概念，如果存在一个多项式的学习算法能够学习它，学习的正确率仅比随机猜测略好，那么就称这个概念是弱可学习的。非常有趣的是 Schapire 后来证明强可学习与弱可学习是等价的，也就是说，在 PAC 学习的框架下，一个概念是强可学习的充分必要条件是这个概念是弱可学习的。

这样一来，问题便成为，在学习中，如果已经发现了“弱学习算法”，那么能否将它提升（boost）为“强学习算法”。大家知道，发现弱学习算法通常要比发现强学习算法容易得多。那么如何具体实施提升，便成为开发提升方法时所要解决的问题。关于提升方法的研究很多，有很多算法被提出。最具代表性的是 AdaBoost 算法（AdaBoost algorithm）。

对于分类问题而言，给定一个训练样本集，求比较粗糙的分类规则（弱分类器）要比求精确的分类规则（强分类器）容易得多。提升方法就是从弱学习算法出发，反复学习，得到一系列弱分类器（又称为基本分类器），然后组合这些弱分类器，构成一个强分类器。大多数的提升方法都是改变训练数据的概率分布（训练数据的权值分布），针对不同的训练数据分布调用弱学习算法学习一系列弱分类器。

这样，对提升方法来说，有两个问题需要回答：一是在每一轮如何改变训练数据的权值或概率分布；二是如何将弱分类器组合成一个强分类器。关于第 1 个问题，AdaBoost 的做法是，提高那些被前一轮弱分类器错误分类样本的权值，而降低那些被正确分类样本的权值。这样一来，那些没有得到正确分类的数据，由于其权值的加大而受到后一轮的弱分类器的更大关注。于是，分类问题被一系列的弱分类器“分而治之”。至于第 2 个问题，即弱分类器的组合，AdaBoost 采取加权多数表决的方法。具体地，加大分类误差率小的弱分类器的权值，使其在表决中起较大的作用；减小分类误差率大的弱分类器的权值，使其在表决中起较小的作用。

AdaBoost 的巧妙之处就在于它将这些想法自然且有效地实现在一种算法里。

\subsection*{8.1.2\quad AdaBoost 算法}

现在叙述 AdaBoost 算法。假设给定一个二类分类的训练数据集
\[
T=\{(x_1,y_1),(x_2,y_2),\cdots,(x_N,y_N)\}
\]
其中，每个样本点由实例与标记组成。实例 $x_i\in\mathcal{X}\subseteq\mathbb{R}^n$，标记 $y_i\in\mathcal{Y}=\{-1,+1\}$，$\mathcal{X}$ 是实例空间，$\mathcal{Y}$ 是标记集合。AdaBoost 利用以下算法，从训练数据中学习一系列弱分类器或基本分类器，并将这些弱分类器线性组合成为一个强分类器。

\begin{center}
\fbox{\begin{minipage}{0.92\textwidth}
\noindent\textbf{算法 8.1（AdaBoost）}

\medskip
\noindent 输入：训练数据集 $T=\{(x_1,y_1),(x_2,y_2),\cdots,(x_N,y_N)\}$，其中 $x_i\in\mathcal{X}\subseteq\mathbb{R}^n$，$y_i\in\mathcal{Y}=\{-1,+1\}$；弱学习算法；

\noindent 输出：最终分类器 $G(x)$。

\medskip
\noindent （1）初始化训练数据的权值分布
\[
D_1=(w_{11},\cdots,w_{1i},\cdots,w_{1N}),\qquad
w_{1i}=\frac{1}{N},\quad i=1,2,\cdots,N.
\]

\noindent （2）对 $m=1,2,\cdots,M$

\quad （a）使用具有权值分布 $D_m$ 的训练数据集学习，得到基本分类器
\[
G_m(x):\mathcal{X}\to\{-1,+1\}.
\]

\quad （b）计算 $G_m(x)$ 在训练数据集上的分类误差率
\begin{equation}
e_m=\sum_{i=1}^N P(G_m(x_i)\ne y_i)
=\sum_{i=1}^N w_{mi}I(G_m(x_i)\ne y_i).
\tag{8.1}
\end{equation}

\quad （c）计算 $G_m(x)$ 的系数
\begin{equation}
\alpha_m=\frac{1}{2}\log\frac{1-e_m}{e_m}.
\tag{8.2}
\end{equation}
这里的对数是自然对数。

\quad （d）更新训练数据集的权值分布
\begin{equation}
D_{m+1}=(w_{m+1,1},\cdots,w_{m+1,i},\cdots,w_{m+1,N})
\tag{8.3}
\end{equation}
\begin{equation}
w_{m+1,i}=\frac{w_{mi}}{Z_m}\exp(-\alpha_m y_iG_m(x_i)),
\quad i=1,2,\cdots,N.
\tag{8.4}
\end{equation}
这里，$Z_m$ 是规范化因子
\begin{equation}
Z_m=\sum_{i=1}^N w_{mi}\exp(-\alpha_m y_iG_m(x_i))
\tag{8.5}
\end{equation}
它使 $D_{m+1}$ 成为一个概率分布。

\noindent （3）构建基本分类器的线性组合
\begin{equation}
f(x)=\sum_{m=1}^M \alpha_mG_m(x)
\tag{8.6}
\end{equation}
得到最终分类器
\begin{equation}
\begin{aligned}
G(x)&=\operatorname{sign}(f(x))\\
&=\operatorname{sign}\left(\sum_{m=1}^M \alpha_mG_m(x)\right).
\end{aligned}
\tag{8.7}
\end{equation}
\end{minipage}}
\end{center}

对 AdaBoost 算法作如下说明：

步骤（1）\quad 假设训练数据集具有均匀的权值分布，即每个训练样本在基本分类器的学习中作用相同，这一假设保证第 1 步能够在原始数据上学习基本分类器 $G_1(x)$。

步骤（2）\quad AdaBoost 反复学习基本分类器，在每一轮 $m=1,2,\cdots,M$ 顺次地执行下列操作：

（a）使用当前分布 $D_m$ 加权的训练数据集，学习基本分类器 $G_m(x)$。

（b）计算基本分类器 $G_m(x)$ 在加权训练数据集上的分类误差率：
\begin{equation}
\begin{aligned}
e_m&=\sum_{i=1}^N P(G_m(x_i)\ne y_i)\\
&=\sum_{G_m(x_i)\ne y_i}w_{mi}.
\end{aligned}
\tag{8.8}
\end{equation}
这里，$w_{mi}$ 表示第 $m$ 轮中第 $i$ 个实例的权值，$\sum_{i=1}^Nw_{mi}=1$。这表明，$G_m(x)$ 在加权的训练数据集上的分类误差率是被 $G_m(x)$ 误分类样本的权值之和，由此可以看出数据权值分布 $D_m$ 与基本分类器 $G_m(x)$ 的分类误差率的关系。

（c）计算基本分类器 $G_m(x)$ 的系数 $\alpha_m$。$\alpha_m$ 表示 $G_m(x)$ 在最终分类器中的重要性。由式（8.2）可知，当 $e_m\leq \frac{1}{2}$ 时，$\alpha_m\geq 0$，并且 $\alpha_m$ 随着 $e_m$ 的减小而增大，所以分类误差率越小的基本分类器在最终分类器中的作用越大。

（d）更新训练数据的权值分布为下一轮作准备。式（8.4）可以写成：
\[
w_{m+1,i}=
\begin{cases}
\dfrac{w_{mi}}{Z_m}e^{-\alpha_m}, & G_m(x_i)=y_i,\\[0.4em]
\dfrac{w_{mi}}{Z_m}e^{\alpha_m}, & G_m(x_i)\ne y_i.
\end{cases}
\]
由此可知，被基本分类器 $G_m(x)$ 误分类样本的权值得以扩大，而被正确分类样本的权值却得以缩小。两相比较，由式（8.2）知误分类样本的权值被放大 $e^{2\alpha_m}=\frac{1-e_m}{e_m}$ 倍。因此，误分类样本在下一轮学习中起更大的作用。不改变所给的训练数据，而不断改变训练数据权值的分布，使得训练数据在基本分类器的学习中起不同的作用，这是 AdaBoost 的一个特点。

步骤（3）\quad 线性组合 $f(x)$ 实现 $M$ 个基本分类器的加权表决。系数 $\alpha_m$ 表示了基本分类器 $G_m(x)$ 的重要性，这里，所有 $\alpha_m$ 之和并不为 1。$f(x)$ 的符号决定实例 $x$ 的类，$f(x)$ 的绝对值表示分类的确信度。利用基本分类器的线性组合构建最终分类器是 AdaBoost 的另一个特点。

\subsection*{8.1.3\quad AdaBoost 的例子\protect\footnote{例题来源于 http://www.csie.edu.tw。}}

\noindent\textbf{例 8.1}\quad 给定如表 8.1 所示训练数据。假设弱分类器由 $x<v$ 或 $x>v$ 产生，其阈值 $v$ 使该分类器在训练数据集上分类误差率最低。试用 AdaBoost 算法学习一个强分类器。

\begin{table}[htbp]
\centering
\caption{训练数据表}
\begin{tabular}{c*{10}{c}}
\toprule
序号 & 1 & 2 & 3 & 4 & 5 & 6 & 7 & 8 & 9 & 10\\
\midrule
$x$ & 0 & 1 & 2 & 3 & 4 & 5 & 6 & 7 & 8 & 9\\
$y$ & 1 & 1 & 1 & $-1$ & $-1$ & $-1$ & 1 & 1 & 1 & $-1$\\
\bottomrule
\end{tabular}
\end{table}

\noindent\textbf{解}\quad 初始化数据权值分布
\[
D_1=(w_{11},w_{12},\cdots,w_{110})
\]
\[
w_{1i}=0.1,\quad i=1,2,\cdots,10.
\]

对 $m=1$，

（a）在权值分布为 $D_1$ 的训练数据上，阈值 $v$ 取 2.5 时分类误差率最低，故基本分类器为
\[
G_1(x)=
\begin{cases}
1, & x<2.5\\
-1, & x>2.5
\end{cases}
\]

（b）$G_1(x)$ 在训练数据集上的误差率 $e_1=P(G_1(x_i)\ne y_i)=0.3$。

（c）计算 $G_1(x)$ 的系数：$\alpha_1=\frac{1}{2}\log\frac{1-e_1}{e_1}=0.4236$。

（d）更新训练数据的权值分布：
\[
D_2=(w_{21},\cdots,w_{2i},\cdots,w_{210})
\]
\[
w_{2i}=\frac{w_{1i}}{Z_1}\exp(-\alpha_1y_iG_1(x_i)),\quad i=1,2,\cdots,10
\]
\[
D_2=(0.07143,0.07143,0.07143,0.07143,0.07143,0.07143,
0.16667,0.16667,0.16667,0.07143)
\]
\[
f_1(x)=0.4236G_1(x)
\]

分类器 $\operatorname{sign}[f_1(x)]$ 在训练数据集上有 3 个误分类点。

对 $m=2$，

（a）在权值分布为 $D_2$ 的训练数据上，阈值 $v$ 是 8.5 时分类误差率最低，基本分类器为
\[
G_2(x)=
\begin{cases}
1, & x<8.5\\
-1, & x>8.5
\end{cases}
\]

（b）$G_2(x)$ 在训练数据集上的误差率 $e_2=0.2143$。

（c）计算 $\alpha_2=0.6496$。

（d）更新训练数据权值分布：
\[
D_3=(0.0455,0.0455,0.0455,0.1667,0.1667,0.1667,
0.1060,0.1060,0.1060,0.0455)
\]
\[
f_2(x)=0.4236G_1(x)+0.6496G_2(x)
\]

分类器 $\operatorname{sign}[f_2(x)]$ 在训练数据集上有 3 个误分类点。

对 $m=3$，

（a）在权值分布为 $D_3$ 的训练数据上，阈值 $v$ 是 5.5 时分类误差率最低，基本分类器为
\[
G_3(x)=
\begin{cases}
1, & x>5.5\\
-1, & x<5.5
\end{cases}
\]

（b）$G_3(x)$ 在训练样本集上的误差率 $e_3=0.1820$。

（c）计算 $\alpha_3=0.7514$。

（d）更新训练数据的权值分布：
\[
D_4=(0.125,0.125,0.125,0.102,0.102,0.102,0.065,0.065,0.065,0.125)
\]
于是得到：
\[
f_3(x)=0.4236G_1(x)+0.6496G_2(x)+0.7514G_3(x)
\]
分类器 $\operatorname{sign}[f_3(x)]$ 在训练数据集上误分类点个数为 0。

于是最终分类器为
\[
G(x)=\operatorname{sign}[f_3(x)]
=\operatorname{sign}[0.4236G_1(x)+0.6496G_2(x)+0.7514G_3(x)].
\]

\subsection*{8.2\quad AdaBoost 算法的训练误差分析}

AdaBoost 最基本的性质是它能在学习过程中不断减少训练误差，即在训练数据集上的分类误差率。关于这个问题有下面的定理。

\noindent\textbf{定理 8.1（AdaBoost 的训练误差界）}\quad
AdaBoost 算法最终分类器的训练误差界为
\begin{equation}
\frac{1}{N}\sum_{i=1}^N I(G(x_i)\ne y_i)
\leq \frac{1}{N}\sum_i \exp(-y_if(x_i))=\prod_m Z_m
\tag{8.9}
\end{equation}
这里，$G(x),f(x)$ 和 $Z_m$ 分别由式（8.7）、式（8.6）和式（8.5）给出。

\noindent\textbf{证明}\quad
当 $G(x_i)\ne y_i$ 时，$y_if(x_i)<0$，因而 $\exp(-y_if(x_i))\geq 1$。由此直接推导出前半部分。

后半部分的推导要用到 $Z_m$ 的定义式（8.5）及式（8.4）的变形：
\[
w_{mi}\exp(-\alpha_my_iG_m(x_i))=Z_mw_{m+1,i}.
\]
现推导如下：
\[
\begin{aligned}
\frac{1}{N}\sum_i \exp(-y_if(x_i))
&=\frac{1}{N}\sum_i \exp\left(-\sum_{m=1}^M \alpha_my_iG_m(x_i)\right)\\
&=\sum_i w_{1i}\prod_{m=1}^M \exp(-\alpha_my_iG_m(x_i))\\
&=Z_1\sum_i w_{2i}\prod_{m=2}^M \exp(-\alpha_my_iG_m(x_i))\\
&=Z_1Z_2\sum_i w_{3i}\prod_{m=3}^M \exp(-\alpha_my_iG_m(x_i))\\
&=\cdots\\
&=Z_1Z_2\cdots Z_{M-1}\sum_i w_{Mi}\exp(-\alpha_My_iG_M(x_i))\\
&=\prod_{m=1}^M Z_m.
\end{aligned}
\]
这一定理说明，可以在每一轮选取适当的 $G_m$ 使得 $Z_m$ 最小，从而使训练误差下降最快。对二类分类问题，有如下结果。

\noindent\textbf{定理 8.2（二类分类问题 AdaBoost 的训练误差界）}\quad
\begin{equation}
\begin{aligned}
\prod_{m=1}^M Z_m
&=\prod_{m=1}^M \left[2\sqrt{e_m(1-e_m)}\right]\\
&=\prod_{m=1}^M \sqrt{1-4\gamma_m^2}\\
&\leq \exp\left(-2\sum_{m=1}^M \gamma_m^2\right)
\end{aligned}
\tag{8.10}
\end{equation}
这里，$\gamma_m=\frac{1}{2}-e_m$。

\noindent\textbf{证明}\quad
由 $Z_m$ 的定义式（8.5）及式（8.8）得
\[
\begin{aligned}
Z_m
&=\sum_{i=1}^N w_{mi}\exp(-\alpha_my_iG_m(x_i))\\
&=\sum_{y_i=G_m(x_i)}w_{mi}e^{-\alpha_m}
+\sum_{y_i\ne G_m(x_i)}w_{mi}e^{\alpha_m}\\
&=(1-e_m)e^{-\alpha_m}+e_me^{\alpha_m}\\
&=2\sqrt{e_m(1-e_m)}\\
&=\sqrt{1-4\gamma_m^2}.
\end{aligned}
\tag{8.11}
\]
至于不等式
\[
\prod_{m=1}^M \sqrt{1-4\gamma_m^2}
\leq \exp\left(-2\sum_{m=1}^M \gamma_m^2\right)
\]
则可先由 $e^x$ 和 $\sqrt{1-x}$ 在点 $x=0$ 的泰勒展开式推出不等式 $\sqrt{1-4\gamma_m^2}\leq \exp(-2\gamma_m^2)$，进而得到。

\noindent\textbf{推论 8.1}\quad
如果存在 $\gamma>0$，对所有 $m$ 有 $\gamma_m\geq \gamma$，则
\begin{equation}
\frac{1}{N}\sum_{i=1}^N I(G(x_i)\ne y_i)\leq \exp(-2M\gamma^2)
\tag{8.12}
\end{equation}

这表明在此条件下 AdaBoost 的训练误差是以指数速率下降的。这一性质当然是很有吸引力的。

注意，AdaBoost 算法不需要知道下界 $\gamma$，这正是 Freund 与 Schapire 设计 AdaBoost 时所考虑的。与一些早期的提升方法不同，AdaBoost 具有适应性，即它能适应弱分类器各自的训练误差率。这也是它的名称（适应的提升）的由来，Ada 是 Adaptive 的简写。

\subsection*{8.3\quad AdaBoost 算法的解释}

AdaBoost 算法还有另一个解释，即可以认为 AdaBoost 算法是模型为加法模型、损失函数为指数函数、学习算法为前向分步算法时的二类分类学习方法。

\subsection*{8.3.1\quad 前向分步算法}

考虑加法模型（additive model）
\begin{equation}
f(x)=\sum_{m=1}^M\beta_m b(x;\gamma_m)
\tag{8.13}
\end{equation}
其中，$b(x;\gamma_m)$ 为基函数，$\gamma_m$ 为基函数的参数，$\beta_m$ 为基函数的系数。显然，式（8.6）是一个加法模型。

在给定训练数据及损失函数 $L(y,f(x))$ 的条件下，学习加法模型 $f(x)$ 成为经验风险极小化即损失函数极小化问题：
\begin{equation}
\min_{\beta_m,\gamma_m}\sum_{i=1}^N
L\left(y_i,\sum_{m=1}^M\beta_m b(x_i;\gamma_m)\right)
\tag{8.14}
\end{equation}

通常这是一个复杂的优化问题。前向分步算法（forward stagewise algorithm）求解这一优化问题的想法是：因为学习的是加法模型，如果能够从前向后，每一步只学习一个基函数及其系数，逐步逼近优化目标函数式（8.14），那么就可以简化优化的复杂度。具体地，每步只需优化如下损失函数：
\begin{equation}
\min_{\beta,\gamma}\sum_{i=1}^N L(y_i,\beta b(x_i;\gamma))
\tag{8.15}
\end{equation}

给定训练数据集 $T=\{(x_1,y_1),(x_2,y_2),\cdots,(x_N,y_N)\}$，$x_i\in\mathcal{X}\subseteq\mathbb{R}^n$，$y_i\in\mathcal{Y}=\{-1,+1\}$。损失函数 $L(y,f(x))$ 和基函数的集合 $\{b(x;\gamma)\}$，学习加法模型 $f(x)$ 的前向分步算法如下。

\begin{center}
\fbox{\begin{minipage}{0.92\textwidth}
\noindent\textbf{算法 8.2（前向分步算法）}

\medskip
\noindent 输入：训练数据集 $T=\{(x_1,y_1),(x_2,y_2),\cdots,(x_N,y_N)\}$；损失函数 $L(y,f(x))$；基函数集 $\{b(x;\gamma)\}$；

\noindent 输出：加法模型 $f(x)$。

\medskip
\noindent （1）初始化 $f_0(x)=0$；

\noindent （2）对 $m=1,2,\cdots,M$

\quad （a）极小化损失函数
\begin{equation}
(\beta_m,\gamma_m)=\arg\min_{\beta,\gamma}
\sum_{i=1}^N L(y_i,f_{m-1}(x_i)+\beta b(x_i;\gamma))
\tag{8.16}
\end{equation}
得到参数 $\beta_m,\gamma_m$。

\quad （b）更新
\begin{equation}
f_m(x)=f_{m-1}(x)+\beta_m b(x;\gamma_m)
\tag{8.17}
\end{equation}

\noindent （3）得到加法模型
\begin{equation}
f(x)=f_M(x)=\sum_{m=1}^M\beta_m b(x;\gamma_m)
\tag{8.18}
\end{equation}
\end{minipage}}
\end{center}

这样，前向分步算法将同时求解从 $m=1$ 到 $M$ 所有参数 $\beta_m,\gamma_m$ 的优化问题简化为逐次求解各个 $\beta_m,\gamma_m$ 的优化问题。

\subsection*{8.3.2\quad 前向分步算法与 AdaBoost}

由前向分步算法可以推导出 AdaBoost，用定理叙述这一关系。

\noindent\textbf{定理 8.3}\quad
AdaBoost 算法是前向分步加法算法的特例。这时，模型是由基本分类器组成的加法模型，损失函数是指数函数。

\noindent\textbf{证明}\quad
前向分步算法学习的是加法模型，当基函数为基本分类器时，该加法模型等价于 AdaBoost 的最终分类器
\begin{equation}
f(x)=\sum_{m=1}^M\alpha_mG_m(x)
\tag{8.19}
\end{equation}
由基本分类器 $G_m(x)$ 及其系数 $\alpha_m$ 组成，$m=1,2,\cdots,M$。前向分步算法逐一学习基函数，这一过程与 AdaBoost 算法逐一学习基本分类器的过程一致。下面证明前向分步算法的损失函数是指数损失函数（exponential loss function）
\[
L(y,f(x))=\exp[-yf(x)]
\]
时，其学习的具体操作等价于 AdaBoost 算法学习的具体操作。

假设经过 $m-1$ 轮迭代前向分步算法已经得到 $f_{m-1}(x)$：
\[
\begin{aligned}
f_{m-1}(x)&=f_{m-2}(x)+\alpha_{m-1}G_{m-1}(x)\\
&=\alpha_1G_1(x)+\cdots+\alpha_{m-1}G_{m-1}(x)
\end{aligned}
\]
在第 $m$ 轮迭代得到 $\alpha_m,G_m(x)$ 和 $f_m(x)$。
\[
f_m(x)=f_{m-1}(x)+\alpha_mG_m(x)
\]
目标是使前向分步算法得到的 $\alpha_m$ 和 $G_m(x)$ 使 $f_m(x)$ 在训练数据集 $T$ 上的指数损失最小，即
\begin{equation}
(\alpha_m,G_m(x))=\arg\min_{\alpha,G}
\sum_{i=1}^N \exp[-y_i(f_{m-1}(x_i)+\alpha G(x_i))]
\tag{8.20}
\end{equation}
式（8.20）可以表示为
\begin{equation}
(\alpha_m,G_m(x))=\arg\min_{\alpha,G}
\sum_{i=1}^N \bar{w}_{mi}\exp[-y_i\alpha G(x_i)]
\tag{8.21}
\end{equation}
其中，$\bar{w}_{mi}=\exp[-y_if_{m-1}(x_i)]$。因为 $\bar{w}_{mi}$ 既不依赖 $\alpha$ 也不依赖于 $G$，所以与最小化无关。但 $\bar{w}_{mi}$ 依赖于 $f_{m-1}(x)$，随着每一轮迭代而发生改变。

现证明使式（8.21）达到最小的 $\alpha_m^\ast$ 和 $G_m^\ast(x)$ 就是 AdaBoost 算法所得到的 $\alpha_m$ 和 $G_m(x)$。求解式（8.21）可分两步：

首先，求 $G_m^\ast(x)$。对任意 $\alpha>0$，使式（8.21）最小的 $G(x)$ 由下式得到：
\[
G_m^\ast(x)=\arg\min_G\sum_{i=1}^N \bar{w}_{mi}I(y_i\ne G(x_i))
\]
其中，$\bar{w}_{mi}=\exp[-y_if_{m-1}(x_i)]$。

此分类器 $G_m^\ast(x)$ 即为 AdaBoost 算法的基本分类器 $G_m(x)$，因为它是使第 $m$ 轮加权训练数据分类误差率最小的基本分类器。

之后，求 $\alpha_m^\ast$。参照式（8.11），式（8.21）中
\begin{equation}
\begin{aligned}
\sum_{i=1}^N\bar{w}_{mi}\exp[-y_i\alpha G(x_i)]
&=\sum_{y_i=G_m(x_i)}\bar{w}_{mi}e^{-\alpha}
+\sum_{y_i\ne G_m(x_i)}\bar{w}_{mi}e^\alpha\\
&=(e^\alpha-e^{-\alpha})\sum_{i=1}^N\bar{w}_{mi}I(y_i\ne G(x_i))
+e^{-\alpha}\sum_{i=1}^N\bar{w}_{mi}.
\end{aligned}
\tag{8.22}
\end{equation}
将已求得的 $G_m^\ast(x)$ 代入式（8.22），对 $\alpha$ 求导并使导数为 0，即得到使式（8.21）最小的 $\alpha$。
\[
\alpha_m^\ast=\frac{1}{2}\log\frac{1-e_m}{e_m}
\]
其中，$e_m$ 是分类误差率：
\begin{equation}
\begin{aligned}
e_m&=\frac{\sum_{i=1}^N\bar{w}_{mi}I(y_i\ne G_m(x_i))}
{\sum_{i=1}^N\bar{w}_{mi}}\\
&=\sum_{i=1}^N w_{mi}I(y_i\ne G_m(x_i)).
\end{aligned}
\tag{8.23}
\end{equation}
这里的 $\alpha_m^\ast$ 与 AdaBoost 算法第 2(c) 步的 $\alpha_m$ 完全一致。

最后来看每一轮样本权值的更新。由
\[
f_m(x)=f_{m-1}(x)+\alpha_mG_m(x)
\]
以及 $\bar{w}_{mi}=\exp[-y_if_{m-1}(x_i)]$，可得
\[
\bar{w}_{m+1,i}=\bar{w}_{mi}\exp[-y_i\alpha_mG_m(x)]
\]
这与 AdaBoost 算法第 2(d) 步的样本权值的更新，只相差规范化因子，因而等价。

\subsection*{8.4\quad 提升树}

提升树是以分类树或回归树为基本分类器的提升方法。提升树被认为是统计学习中性能最好的方法之一。

\subsection*{8.4.1\quad 提升树模型}

提升方法实际采用加法模型（即基函数的线性组合）与前向分步算法。以决策树为基函数的提升方法称为提升树（boosting tree）。对分类问题决策树是二叉分类树，对回归问题决策树是二叉回归树。在例 8.1 中看到的基本分类器 $x<v$ 或 $x>v$，可以看作是由一个根结点直接连接两个叶结点的简单决策树，即所谓的决策树桩（decision stump）。提升树模型可以表示为决策树的加法模型：
\begin{equation}
f_M(x)=\sum_{m=1}^M T(x;\Theta_m)
\tag{8.24}
\end{equation}
其中，$T(x;\Theta_m)$ 表示决策树，$\Theta_m$ 为决策树的参数，$M$ 为树的个数。

\subsection*{8.4.2\quad 提升树算法}

提升树算法采用前向分步算法。首先确定初始提升树 $f_0(x)=0$，第 $m$ 步的模型是
\begin{equation}
f_m(x)=f_{m-1}(x)+T(x;\Theta_m)
\tag{8.25}
\end{equation}
其中，$f_{m-1}(x)$ 为当前模型，通过经验风险极小化确定下一棵决策树的参数 $\Theta_m$：
\begin{equation}
\hat{\Theta}_m=\arg\min_{\Theta_m}\sum_{i=1}^N
L(y_i,f_{m-1}(x_i)+T(x_i;\Theta_m))
\tag{8.26}
\end{equation}

由于树的线性组合可以很好地拟合训练数据，即使数据中的输入与输出之间的关系很复杂也是如此，所以提升树是一个高功能的学习算法。

下面讨论针对不同问题的提升树学习算法，其主要区别在于使用的损失函数不同。包括用平方误差损失函数的回归问题，用指数损失函数的分类问题，以及用一般损失函数的一般决策问题。

对于二类分类问题，提升树算法只需将 AdaBoost 算法 8.1 中的基本分类器限制为二类分类树即可，可以说这时的提升树算法是 AdaBoost 算法的特殊情况，这里不再细述。下面叙述回归问题的提升树。

已知一个训练数据集 $T=\{(x_1,y_1),(x_2,y_2),\cdots,(x_N,y_N)\}$，$x_i\in\mathcal{X}\subseteq\mathbb{R}^n$，$\mathcal{X}$ 为输入空间，$y_i\in\mathcal{Y}\subseteq\mathbb{R}$，$\mathcal{Y}$ 为输出空间。在 5.5 节中已经讨论了回归树的问题。如果将输入空间 $\mathcal{X}$ 划分为 $J$ 个互不相交的区域 $R_1,R_2,\cdots,R_J$，并且在每个区域上确定输出的常量 $c_j$，那么树可表示为
\begin{equation}
T(x;\Theta)=\sum_{j=1}^J c_jI(x\in R_j)
\tag{8.27}
\end{equation}
其中，参数 $\Theta=\{(R_1,c_1),(R_2,c_2),\cdots,(R_J,c_J)\}$ 表示树的区域划分和各区域上的常数。$J$ 是回归树的复杂度即叶结点个数。

回归问题提升树使用以下前向分步算法：
\[
f_0(x)=0
\]
\[
f_m(x)=f_{m-1}(x)+T(x;\Theta_m),\quad m=1,2,\cdots,M
\]
\[
f_M(x)=\sum_{m=1}^M T(x;\Theta_m)
\]
在前向分步算法的第 $m$ 步，给定当前模型 $f_{m-1}(x)$，需求解
\[
\hat{\Theta}_m=\arg\min_{\Theta_m}\sum_{i=1}^N
L(y_i,f_{m-1}(x_i)+T(x_i;\Theta_m))
\]
得到 $\hat{\Theta}_m$，即第 $m$ 棵树的参数。

当采用平方误差损失函数时，
\[
L(y,f(x))=(y-f(x))^2
\]
其损失变为
\[
\begin{aligned}
L(y,f_{m-1}(x)+T(x;\Theta_m))
&=[y-f_{m-1}(x)-T(x;\Theta_m)]^2\\
&=[r-T(x;\Theta_m)]^2
\end{aligned}
\]
这里，
\begin{equation}
r=y-f_{m-1}(x)
\tag{8.28}
\end{equation}
是当前模型拟合数据的残差（residual）。所以，对回归问题的提升树算法来说，只需简单地拟合当前模型的残差。这样，算法是相当简单的。现将回归问题的提升树算法叙述如下。

\begin{center}
\fbox{\begin{minipage}{0.92\textwidth}
\noindent\textbf{算法 8.3（回归问题的提升树算法）}

\medskip
\noindent 输入：训练数据集 $T=\{(x_1,y_1),(x_2,y_2),\cdots,(x_N,y_N)\}$，$x_i\in\mathcal{X}\subseteq\mathbb{R}^n$，$y_i\in\mathcal{Y}\subseteq\mathbb{R}$；

\noindent 输出：提升树 $f_M(x)$。

\medskip
\noindent （1）初始化 $f_0(x)=0$。

\noindent （2）对 $m=1,2,\cdots,M$。

\quad （a）按式（8.27）计算残差：
\[
r_{mi}=y_i-f_{m-1}(x_i),\quad i=1,2,\cdots,N
\]

\quad （b）拟合残差 $r_{mi}$ 学习一个回归树，得到 $T(x;\Theta_m)$。

\quad （c）更新 $f_m(x)=f_{m-1}(x)+T(x;\Theta_m)$。

\noindent （3）得到回归问题提升树
\[
f_M(x)=\sum_{m=1}^M T(x;\Theta_m)
\]
\end{minipage}}
\end{center}

\noindent\textbf{例 8.2}\quad
已知如表 8.2 所示的训练数据，$x$ 的取值范围为区间 $[0.5,10.5]$，$y$ 的取值范围为区间 $[5.0,10.0]$，学习这个回归问题的提升树模型，考虑只用树桩作为基函数。

\begin{table}[htbp]
\centering
\caption{训练数据表}
\begin{tabular}{c*{10}{c}}
\toprule
$x_i$ & 1 & 2 & 3 & 4 & 5 & 6 & 7 & 8 & 9 & 10\\
\midrule
$y_i$ & 5.56 & 5.70 & 5.91 & 6.40 & 6.80 & 7.05 & 8.90 & 8.70 & 9.00 & 9.05\\
\bottomrule
\end{tabular}
\end{table}

\noindent\textbf{解}\quad
按照算法 8.3，第 1 步求 $f_1(x)$ 即回归树 $T_1(x)$。

首先通过以下优化问题：
\[
\min_s\left[
\min_{c_1}\sum_{x_i\in R_1}(y_i-c_1)^2
+\min_{c_2}\sum_{x_i\in R_2}(y_i-c_2)^2
\right]
\]
求解训练数据的切分点 $s$：
\[
R_1=\{x|x\leq s\},\qquad R_2=\{x|x>s\}
\]
容易求得使 $R_1,R_2$ 内部使平方损失误差达到最小值的 $c_1,c_2$ 为
\[
c_1=\frac{1}{N_1}\sum_{x_i\in R_1}y_i,\qquad
c_2=\frac{1}{N_2}\sum_{x_i\in R_2}y_i
\]
这里 $N_1,N_2$ 是 $R_1,R_2$ 的样本点数。

求训练数据的切分点。根据所给数据，考虑如下切分点：
\[
1.5,\ 2.5,\ 3.5,\ 4.5,\ 5.5,\ 6.5,\ 7.5,\ 8.5,\ 9.5
\]

对各切分点，不难求出相应的 $R_1,R_2,c_1,c_2$ 及
\[
m(s)=\min_{c_1}\sum_{x_i\in R_1}(y_i-c_1)^2
+\min_{c_2}\sum_{x_i\in R_2}(y_i-c_2)^2.
\]
例如，当 $s=1.5$ 时，$R_1=\{1\}$，$R_2=\{2,3,\cdots,10\}$，$c_1=5.56$，$c_2=7.50$，
\[
m(s)=\min_{c_1}\sum_{x_i\in R_1}(y_i-c_1)^2
+\min_{c_2}\sum_{x_i\in R_2}(y_i-c_2)^2=0+15.72=15.72
\]

现将 $s$ 及 $m(s)$ 的计算结果列表如下（见表 8.3）。

\begin{table}[htbp]
\centering
\caption{计算数据表}
\begin{tabular}{c*{9}{c}}
\toprule
$s$ & 1.5 & 2.5 & 3.5 & 4.5 & 5.5 & 6.5 & 7.5 & 8.5 & 9.5\\
\midrule
$m(s)$ & 15.72 & 12.07 & 8.36 & 5.78 & 3.91 & 1.93 & 8.01 & 11.73 & 15.74\\
\bottomrule
\end{tabular}
\end{table}

由表 8.3 可知，当 $s=6.5$ 时 $m(s)$ 达到最小值，此时 $R_1=\{1,2,\cdots,6\}$，$R_2=\{7,8,9,10\}$，$c_1=6.24,c_2=8.91$，所以回归树 $T_1(x)$ 为
\[
T_1(x)=
\begin{cases}
6.24, & x<6.5\\
8.91, & x\geq 6.5
\end{cases}
\]
\[
f_1(x)=T_1(x)
\]

用 $f_1(x)$ 拟合训练数据的残差见表 8.4，表中 $r_{2i}=y_i-f_1(x_i)$，$i=1,2,\cdots,10$。

\begin{table}[htbp]
\centering
\caption{残差表}
\begin{tabular}{c*{10}{c}}
\toprule
$x_i$ & 1 & 2 & 3 & 4 & 5 & 6 & 7 & 8 & 9 & 10\\
\midrule
$r_{2i}$ & $-0.68$ & $-0.54$ & $-0.33$ & 0.16 & 0.56 & 0.81 & $-0.01$ & $-0.21$ & 0.09 & 0.14\\
\bottomrule
\end{tabular}
\end{table}

用 $f_1(x)$ 拟合训练数据的平方损失误差：
\[
L(y,f_1(x))=\sum_{i=1}^{10}(y_i-f_1(x_i))^2=1.93
\]

第 2 步求 $T_2(x)$。方法与求 $T_1(x)$ 一样，只是拟合的数据是表 8.4 的残差。可以得到：
\[
T_2(x)=
\begin{cases}
-0.52, & x<3.5\\
0.22, & x\geq 3.5
\end{cases}
\]
\[
f_2(x)=f_1(x)+T_2(x)=
\begin{cases}
5.72, & x<3.5\\
6.46, & 3.5\leq x<6.5\\
9.13, & x\geq 6.5
\end{cases}
\]

用 $f_2(x)$ 拟合训练数据的平方损失误差是
\[
L(y,f_2(x))=\sum_{i=1}^{10}(y_i-f_2(x_i))^2=0.79
\]

继续求得
\[
T_3(x)=
\begin{cases}
0.15, & x<6.5\\
-0.22, & x\geq 6.5
\end{cases}
\qquad
L(y,f_3(x))=0.47,
\]
\[
T_4(x)=
\begin{cases}
-0.16, & x<4.5\\
0.11, & x\geq 4.5
\end{cases}
\qquad
L(y,f_4(x))=0.30,
\]
\[
T_5(x)=
\begin{cases}
0.07, & x<6.5\\
-0.11, & x\geq 6.5
\end{cases}
\qquad
L(y,f_5(x))=0.23,
\]
\[
T_6(x)=
\begin{cases}
-0.15, & x<2.5\\
0.04, & x\geq 2.5
\end{cases}
\]
\[
\begin{aligned}
f_6(x)&=f_5(x)+T_6(x)=T_1(x)+\cdots+T_5(x)+T_6(x)\\
&=
\begin{cases}
5.63, & x<2.5\\
5.82, & 2.5\leq x<3.5\\
6.56, & 3.5\leq x<4.5\\
6.83, & 4.5\leq x<6.5\\
8.95, & x\geq 6.5
\end{cases}
\end{aligned}
\]
用 $f_6(x)$ 拟合训练数据的平方损失误差是
\[
L(y,f_6(x))=\sum_{i=1}^{10}(y_i-f_6(x_i))^2=0.17
\]
假设此时已满足误差要求，那么 $f(x)=f_6(x)$ 即为所求提升树。

\subsection*{8.4.3\quad 梯度提升}

提升树利用加法模型与前向分步算法实现学习的优化过程。当损失函数是平方损失和指数损失函数时，每一步优化是很简单的。但对一般损失函数而言，往往每一步优化并不那么容易。针对这一问题，Freidman 提出了梯度提升（gradient boosting）算法。这是利用最速下降法的近似方法，其关键是利用损失函数的负梯度在当前模型的值
\[
-\left[\frac{\partial L(y,f(x_i))}{\partial f(x_i)}\right]_{f(x)=f_{m-1}(x)}
\]
作为回归问题提升树算法中的残差的近似值，拟合一个回归树。

\begin{center}
\fbox{\begin{minipage}{0.92\textwidth}
\noindent\textbf{算法 8.4（梯度提升算法）}

\medskip
\noindent 输入：训练数据集 $T=\{(x_1,y_1),(x_2,y_2),\cdots,(x_N,y_N)\}$，$x_i\in\mathcal{X}\subseteq\mathbb{R}^n$，$y_i\in\mathcal{Y}\subseteq\mathbb{R}$；损失函数 $L(y,f(x))$；

\noindent 输出：回归树 $\hat{f}(x)$。

\medskip
\noindent （1）初始化
\[
f_0(x)=\arg\min_c\sum_{i=1}^N L(y_i,c)
\]

\noindent （2）对 $m=1,2,\cdots,M$

\quad （a）对 $i=1,2,\cdots,N$，计算
\[
r_{mi}=-\left[\frac{\partial L(y_i,f(x_i))}{\partial f(x_i)}\right]_{f(x)=f_{m-1}(x)}
\]

\quad （b）对 $r_{mi}$ 拟合一个回归树，得到第 $m$ 棵树的叶结点区域 $R_{mj}$，$j=1,2,\cdots,J$。

\quad （c）对 $j=1,2,\cdots,J$，计算
\[
c_{mj}=\arg\min_c\sum_{x_i\in R_{mj}}L(y_i,f_{m-1}(x_i)+c)
\]

\quad （d）更新
\[
f_m(x)=f_{m-1}(x)+\sum_{j=1}^J c_{mj}I(x\in R_{mj})
\]

\noindent （3）得到回归树
\[
\hat{f}(x)=f_M(x)=\sum_{m=1}^M\sum_{j=1}^J c_{mj}I(x\in R_{mj})
\]
\end{minipage}}
\end{center}

算法第 1 步初始化，估计使损失函数极小化的常数值，它是只有一个根结点的树。第 2(a) 步计算损失函数的负梯度在当前模型的值，将它作为残差的估计。对于平方损失函数，它就是通常所说的残差；对于一般损失函数，它就是残差的近似值。第 2(b) 步估计回归树叶结点区域，以拟合残差的近似值。第 2(c) 步利用线性搜索估计叶结点区域的值，使损失函数极小化。第 2(d) 步更新回归树。第 3 步得到输出的最终模型 $\hat{f}(x)$。

\section*{本章概要}

1. 提升方法是将弱学习算法提升为强学习算法的统计学习方法。在分类学习中，提升方法通过反复修改训练数据的权值分布，构建一系列基本分类器（弱分类器），并将这些基本分类器线性组合，构成一个强分类器。代表性的提升方法是 AdaBoost 算法。

AdaBoost 模型是弱分类器的线性组合：
\[
f(x)=\sum_{m=1}^M \alpha_mG_m(x)
\]

2. AdaBoost 算法的特点是通过迭代每次学习一个基本分类器。每次迭代中，提高那些被前一轮分类器错误分类数据的权值，而降低那些被正确分类的数据的权值。最后，AdaBoost 将基本分类器的线性组合作为强分类器，其中给分类误差率小的基本分类器以大的权值，给分类误差率大的基本分类器以小的权值。

3. AdaBoost 的训练误差分析表明，AdaBoost 的每次迭代可以减少它在训练数据集上的分类误差率，这说明了它作为提升方法的有效性。

4. AdaBoost 算法的一个解释是该算法实际是前向分步算法的一个实现。在这个方法里，模型是加法模型，损失函数是指数损失，算法是前向分步算法。

每一步中极小化损失函数
\[
(\beta_m,\gamma_m)=\arg\min_{\beta,\gamma}\sum_{i=1}^N
L(y_i,f_{m-1}(x_i)+\beta b(x_i;\gamma))
\]
得到参数 $\beta_m,\gamma_m$。

5. 提升树是以分类树或回归树为基本分类器的提升方法。提升树被认为是统计学习中最有效的方法之一。

\section*{继续阅读}

提升方法的介绍可参见文献 [1, 2]。PAC 学习可参见文献 [3]。强可学习与弱可学习的关系可参见文献 [4]。关于 AdaBoost 的最初论文是文献 [5]。关于 AdaBoost 的前向分步加法模型解释参见文献 [6]，提升树与梯度提升可参见文献 [6, 7]。AdaBoost 只是用于二类分类，Schapire 与 Singer 将它扩展到多类分类问题 [8]。AdaBoost 与逻辑斯谛回归的关系也有相关研究 [9]。

\section*{习题}

\noindent\textbf{8.1}\quad
某公司招聘职员考查身体、业务能力、发展潜力这 3 项。身体分为合格 1、不合格 0 两级，业务能力和发展潜力分为上 1、中 2、下 3 三级。分类为合格 1、不合格 $-1$ 两类。已知 10 个人的数据，如表 8.5 所示。假设弱分类器为决策树桩。试用 AdaBoost 算法学习一个强分类器。

\begin{table}[htbp]
\centering
\caption{应聘人员情况数据表}
\begin{tabular}{c*{10}{c}}
\toprule
 & 1 & 2 & 3 & 4 & 5 & 6 & 7 & 8 & 9 & 10\\
\midrule
身体 & 0 & 0 & 1 & 1 & 1 & 0 & 1 & 1 & 1 & 0\\
业务能力 & 1 & 3 & 2 & 1 & 2 & 1 & 1 & 1 & 3 & 2\\
发展潜力 & 3 & 1 & 2 & 3 & 3 & 2 & 2 & 1 & 1 & 1\\
分类 & $-1$ & $-1$ & $-1$ & $-1$ & $-1$ & $-1$ & 1 & 1 & $-1$ & $-1$\\
\bottomrule
\end{tabular}
\end{table}

\noindent\textbf{8.2}\quad
比较支持向量机、AdaBoost、逻辑斯谛回归模型的学习策略与算法。

\section*{参考文献}

\begin{enumerate}[label={[\arabic*]},leftmargin=2.8em]
\item Freund Y, Schapire R E. A short introduction to boosting. Journal of Japanese Society for Artificial Intelligence, 1999, 14(5): 771--780.
\item Hastie T, Tibshirani R, Friedman J. The elements of statistical learning: data mining, inference, and prediction. Springer-Verlag, 2001.（中译本：统计学习基础——数据挖掘、推理与预测. 范明，柴玉梅，眭红英等译. 北京：电子工业出版社，2004.）
\item Valiant L G. A theory of the learnable. Communications of the ACM, 1984, 27(11): 1134--1142.
\item Schapire R. The strength of weak learnability. Machine Learning, 1990, 5(2): 197--227.
\item Freund Y, Schapire R E. A decision-theoretic generalization of on-line learning and an application to boosting. Computational Learning Theory. Lecture Notes in Computer Science, Vol. 904, 1995, 23--37.
\item Friedman J, Hastie T, Tibshirani R. Additive logistic regression: a statistical view of boosting (with discussions). Annals of Statistics, 2000, 28: 337--407.
\item Friedman J. Greedy function approximation: a gradient boosting machine. Annals of Statistics, 2001, 29(5): 1189--1232.
\item Schapire R E, Singer Y. Improved boosting algorithms using confidence-rated predictions. Machine Learning, 1999, 37(3): 297--336.
\item Collins M, Schapire R E, Singer Y. Logistic regression, AdaBoost and Bregman distances. Machine Learning, 2002, 48(1--3): 253--285.
\end{enumerate}

\end{document}
